\section{Introduction}

A zero-knowledge proof can establish that a computation was executed correctly
without revealing all of its inputs.  For a game, this suggests a statement of
the form
\[
  S_{t+1} = F(S_t, I_t),
\]
where $S_t$ is gameplay state, $I_t$ is a private frame input, and $F$ is the
game's transition function.  A proof could commit to the game version, initial
state, route, input stream, and terminal claim while keeping the frame inputs
private.

The equation hides the hard part: what exactly is $F$?  \thgame{} is a closed,
legacy executable whose reconstructed source mixes gameplay with animation,
effects, message scripts, floating-point library behavior, and global random
number consumption.  A deterministic rewrite is not automatically equivalent
to the executable.  Proving the rewrite would only make its differences
cryptographically undeniable.

\project{} therefore places differential semantics before proof engineering.
The immediate objective is a replay runner that can expose the first differing
field at the first differing frame between an instrumented original or exact
reference and a portable implementation.  Only after representative full
replays agree do we freeze a proof-oriented kernel and select a zkVM backend.

Our intended claim is deliberately narrow.  It proves that an input sequence is
a valid execution under a committed game transition and reaches an explicit
terminal predicate.  It does not prove that a human generated the inputs, that
no external tool was used, or that a replay file is authentic.
